\section{Consistency-based diagnosis: proofs}
\label{sec:annex-p16-consistency-based-diagnosis}

Use the sentences, structures and finite component universe of
\cref{sec:p16-consistency-based-diagnosis}. The diagnosis and
measurement arguments follow the formal definitions of
\cite{Reiter1987Diagnosis}. The search proof supplies the
shared-node and relabeling premises explicitly, following
\cite[sec.~4, author version]{GreinerSmithWilkerson1989Correction}.

\begin{annexitem}{annex:p16.diag-semantics}{thm:p16-diag-semantics}{Withdrawal and exact abnormality}
\begin{proof}
For a structure \(\mathcal M\models W^-_\Delta\), put
\[
             A(\mathcal M)=\{c\in C:\mathcal M\models a_c\}.
\]
Bivalence implies \(A(\mathcal M)\subseteq\Delta\), and
\(\mathcal M\models W^=_{A(\mathcal M)}\).
Consequently every consistent withdrawal has a consistent exact
assignment contained in it.

Suppose first that \(\Delta\) is an inclusion-minimal exact
assignment. It is a withdrawal set. If a proper subset
\(D\subsetneq\Delta\) were a withdrawal set, its witnessing
structure would yield an exact assignment
\(A(\mathcal M)\subseteq D\), contrary to minimality of
\(\Delta\). Thus \(\Delta\) is a minimal withdrawal.
Conversely, if \(\Delta\) is a minimal withdrawal, every
structure of \(W^-_\Delta\) has \(A(\mathcal M)=\Delta\);
otherwise it witnesses a smaller withdrawal. Hence
\(W^-_\Delta\entails a_c\) for all \(c\in\Delta\),
and its structures are exactly those of \(W^=_\Delta\).
A smaller exact assignment would be a smaller withdrawal,
so \(\Delta\) is a diagnosis.

Removing normality assumptions preserves a witnessing structure.
Withdrawal sets are therefore upward closed. In a finite
component universe, every nonempty family of withdrawal subsets
of a fixed withdrawal set has an inclusion-minimal member.
That member is a diagnosis by the equivalence just proved.
If \(W\) is consistent, \(C\) is a withdrawal set and thus
contains a diagnosis. If a diagnosis exists, its structure
satisfies \(W\). This proves the existence criterion, including
inconsistent input.

The empty set is a withdrawal exactly when all normality
assumptions are consistent with \(W\). It is then minimal,
and every nonempty withdrawal contains it and is nonminimal.
For \(C=\varnothing\) this says that a consistent \(W\)
has the one empty diagnosis and an inconsistent \(W\)
has none.

Finally, if \(D\subsetneq\Delta\) is a withdrawal set, choose
\(c\in\Delta\setminus D\). Then
\(D\subseteq\Delta\setminus\{c\}\), so that single deletion
is also a withdrawal. This proves the deletion criterion.
For \(W=\{a_c\to a_d\}\), the all-normal assignment is a
model, while exactly \(c\) abnormal violates the implication.
The assignment with both normal still witnesses withdrawal
of \(c\). Thus upward closure of exact assignments fails.
\end{proof}
\end{annexitem}

\begin{annexitem}{annex:p16.diag-conflicts}{thm:p16-diag-conflicts}{Conflict and transversal duality}
\begin{proof}
Adding normality assumptions cannot give a model to an
inconsistent theory; conflicts are upward closed.
If \(W^-_\Delta\) is consistent and a conflict \(K\)
were disjoint from \(\Delta\), then
\(K\subseteq C\setminus\Delta\), and the structure of
\(W^-_\Delta\) would satisfy
\(W\cup\mathsf{Normal}(K)\), a contradiction.
Conversely, if \(W^-_\Delta\) is inconsistent, then
\(C\setminus\Delta\) is a conflict which \(\Delta\)
does not hit. This proves the first equivalence.
Taking minimal sets proves the first diagnosis identity.
Every conflict contains a minimal conflict, by finiteness.
Hitting all minimal conflicts therefore hits all conflicts,
and the second identity follows.

If \(H\) is a minimal hitting set and \(c\in H\), then
\(H\setminus\{c\}\) fails to hit some \(K_c\).
Since \(H\) does hit it, \(K_c\cap H=\{c\}\).
Conversely, these witness conflicts show that every proper
subset of \(H\) misses a conflict: choose an omitted \(c\)
and use \(K_c\).

We prove the double-transversal identity for an arbitrary
finite family \(\mathcal F\). A set \(A\subseteq C\)
hits every member of \(\mathsf{Tr}(\mathcal F)\) exactly
when \(A\) contains some member of \(\mathcal F\).
Indeed, if \(K\in\mathcal F\) and \(K\subseteq A\),
every transversal intersects \(K\), hence \(A\).
If no member of \(\mathcal F\) is contained in \(A\),
then \(C\setminus A\) intersects every member of
\(\mathcal F\). It contains a minimal transversal
disjoint from \(A\). Taking minimal such \(A\) proves
\(\mathsf{Tr}(\mathsf{Tr}(\mathcal F))
=\min_{\subseteq}\mathcal F\).
This reasoning covers \(\mathcal F=\varnothing\):
its minimal transversal is \(\varnothing\), which
no \(A\) hits. If \(\varnothing\in\mathcal F\),
there is no transversal and every \(A\) satisfies the
vacuous second hitting condition, whose minimal member
is \(\varnothing\).
\end{proof}
\end{annexitem}

\begin{annexitem}{annex:p16.diag-singletons}{thm:p16-diag-singletons}{Single faults from conflicts}
\begin{proof}
When a conflict exists, \(\varnothing\) is not a hitting
set. Thus a singleton is a minimal hitting set exactly
when it intersects every conflict, which says precisely
that its element belongs to their intersection.
The conflict-minimal-conflict equivalence just proved
gives the corresponding minimal-conflict intersection.
An empty conflict makes both intersections empty and
correctly gives no singleton.

Every diagnosis hits each known \(K_i\); this proves
the necessary intersection condition. If \(c\) is in
that intersection and \(W^-_{\{c\}}\) is consistent,
then \(\{c\}\) is a withdrawal. It is minimal because
the existence of any \(K_i\) rules out the empty
withdrawal. The withdrawal characterization completes
the converse. Without a conflict, all normal assumptions
are consistent, so the empty diagnosis is the only one.
\end{proof}
\end{annexitem}

\begin{annexitem}{annex:p16.diag-hs-tree}{thm:p16-diag-hs-tree}{The basic hitting-set construction}
\begin{proof}
A conflict used at \(H\) is disjoint from \(H\).
Every outgoing edge therefore adds a new element;
the path length equals \(|H|\), at most \(|C|\).
There are finitely many nodes because the tree has
finite depth and at most \(|C|\) children at a node.
A success label means that no member of \(\mathcal F\)
is disjoint from \(H\), exactly the hitting condition.

Let \(D\) be a minimal hitting set. Starting at the root,
maintain \(H\subseteq D\). At a conflict label \(K\),
choose \(c\in K\cap D\), which exists because \(D\)
hits \(K\). Its child satisfies
\(H\cup\{c\}\subseteq D\). This strictly increases
cardinality, so the process reaches a success node.
It cannot reach an empty conflict, since \(D\) hits
every conflict. At success \(H\subseteq D\) is itself
a hitting set, so minimality gives \(H=D\).

Every nonminimal successful set contains a minimal
hitting set, which also occurs among the successful
sets by the preceding argument. Removing nonminimal
success sets consequently returns exactly the required
minimal family. The proof used only membership and
disjointness of the chosen labels, so arbitrary conflict
orders and nonminimal conflicts are allowed.
\end{proof}
\end{annexitem}

\begin{annexitem}{annex:p16.diag-hs-dag}{thm:p16-diag-hs-dag}{Shared nodes, relabeling and complete pruning}
\begin{proof}
Write \(n=|C|\). Store the set \(H\) as a node's identity.
An edge from \(H\) labeled \(c\) has endpoint
\(H\cup\{c\}\), with \(c\notin H\). Its endpoint
has rank \(|H|+1\). Thus the graph is acyclic, every
root path to \(H\) has length \(|H|\), and sharing
nodes cannot change their represented sets.

If \(K\subsetneq K'\) are in the family, deleting
\(K'\) preserves all hitting sets, not only the minimal
ones: every set hitting \(K\) also hits \(K'\).
At a node with path set \(H\) and old label \(K'\),
one has \(K'\cap H=\varnothing\), hence
\(K\cap H=\varnothing\). Before relabeling there is
an outgoing edge for every element of \(K'\).
Keeping exactly those for \(K\) therefore leaves
every required successor of the new label.
Removing unreachable nodes cannot remove an outgoing
edge from a still reachable node to a required successor:
that edge would itself make its endpoint reachable.
In particular another incoming path preserves the
shared endpoint and all its reachable descendants.

At each completed processing step the following properties
hold for the reachable graph. Each expanded conflict node
has a valid conflict disjoint from its path set and an
edge for every member of that conflict. Every success
node is a hitting set for the original family. A closed
node has a certified successful proper subset as its
closing witness. Unprocessed nodes have not yet been
assigned a completed expansion. The preceding paragraph
shows that relabeling and garbage removal preserve these
properties. The validity of an earlier success certificate
does not depend on retaining its drawing in the graph:
it is an already established hitting set.

At termination, fix a minimal hitting set \(D\).
From the root follow the same choices as in the tree
proof, always maintaining \(H\subseteq D\).
Every conflict label intersects \(D\) and has the needed
outgoing edge. A closed node \(H\subseteq D\) is
impossible: its witness \(J\subsetneq H\subseteq D\)
would be a smaller hitting set than \(D\).
There is no unprocessed reachable node at termination.
The finite increasing path therefore reaches a success
node, and minimality forces its set to be \(D\).
Thus pruning loses no minimal answer.

To show that a success set \(B\) is minimal when first
marked, suppose it has a smaller minimal hitting subset
\(D\subsetneq B\). At that moment every reachable
unprocessed node of rank less than \(|B|\) has already
been processed, by the breadth-first rule. Following
the preceding \(D\)-guided path in the current graph
therefore reaches a successful node with set \(D\);
it cannot end at an unprocessed node, since the entire
path has rank at most \(|D|<|B|\).
The closing rule would then close \(B\), rather than
mark it successful. This contradiction proves minimality.
Relabeling adds no new lower-rank edges or nodes; it only
retains existing edges. Processing the smallest currently
unprocessed rank therefore preserves this breadth-first
argument after a pruning event. Answers occur in
nondecreasing cardinality.

There are at most \(2^n\) distinct active identities and
at most \(n\) outgoing edges per node. For termination,
let \(q\) be the initial number of listed conflicts.
Each nontrivial pruning removes at least one member
from the current family, so there are at most \(q\)
such events. Between them nodes are added and processed
only once for each identity. Even allowing recreation
of identities after a pruning event, at most
\((q+1)2^n\) creations occur. Every processing and
relabeling operation is finite. Hence the procedure
terminates with exactly the asserted answers.

The empty-family root is immediately successful.
If the family contains the empty conflict, any
disjoint-conflict search must eventually expose the
absence of hitting sets; relabeling by that conflict
leaves a root with no children and no success.
These cases also satisfy the proof.
\end{proof}
\end{annexitem}

\begin{annexitem}{annex:p16.diag-oracle}{thm:p16-diag-oracle}{On-demand conflicts and decidable inputs}
\begin{proof}
Put \(n=|C|\).
At a path set \(H\), a cached conflict disjoint from
\(H\) is a valid expansion label. If there is none,
the oracle on \(C\setminus H\) either supplies just
such a conflict or certifies consistency of
\(W^-_H\). By conflict duality, the latter certificate
means that \(H\) hits every actual conflict, including
those never explicitly computed. Thus every expanded,
closed and successful node has exactly the invariant
used in the preceding graph proof, now relative to
the whole implicit family \(\mathcal K(W)\).
That proof of completeness and minimality applies
without enumerating the family.

For every conflict \(K\) previously returned, the cache
always contains some \(K^*\subseteq K\): it initially
contains \(K\), and each redundancy deletion replaces
it by a retained smaller conflict. Suppose a later oracle
call returned \(L\supseteq K\) at \(H\). Since an
oracle conflict is disjoint from \(H\), so would be
\(K^*\subseteq K\subseteq L\). The reuse rule would
have selected a cached conflict and prevented that
oracle call. This is a contradiction.
Consequently returned conflicts are distinct, with at
most \(2^n\) returns. Each pruning needs a new smaller
returned conflict. There are therefore only finitely
many pruning events. Between events the finite
set-identity graph is processed as above, and each
oracle call terminates by hypothesis. The whole
procedure terminates.

For the truncated search, fix a diagnosis \(D\) of
size at most \(k\). The \(D\)-guided path uses ranks
at most \(|D|\), so processing all reachable ranks
at most \(k\) finds it. Successful sets at those ranks
are already minimal by the breadth-first argument.
A later expansion cannot reveal a previously missing
diagnosis of size at most \(k\).

For finite propositional input, evaluate every formula
under each of the \(2^v\) valuations. A satisfying
valuation certifies consistency. If none exists,
the entire queried normality set \(N\) is a conflict
by definition, so a conflict is available without
extracting a minimal proof support. To minimize it,
try each element once: delete it if inconsistency
persists. If deletion of \(c\) is rejected at some
stage, the remaining assumptions without \(c\)
have a model. Subsequent deletions only remove
assumptions, so that model still witnesses consistency
without \(c\) at the end. Every final element is
therefore indispensable. At most \(|N|\) additional
tests are used.

An arbitrary first-order refutation enumeration can
certify inconsistency by terminating with a refutation,
but failure to find one after a finite search is no
consistency certificate. The oracle condition explicitly
requires both kinds of answer and termination on each
query; none is inferred from a merely sound refutation
routine.
\end{proof}
\end{annexitem}

\begin{annexitem}{annex:p16.diag-pruning-calculations}{thm:p16-diag-pruning-calculations}{Three complete set calculations}
\begin{proof}
In the first family every hitting set contains \(2\)
or \(4\), because it must hit \(\{2,4\}\).
If it contains \(2\), the only further requirements
not already met are \(\{1,3,5\}\) and \(\{1,6\}\).
Their minimal simultaneous choices are \(\{1\}\),
\(\{3,6\}\) and \(\{5,6\}\). These give the first
three sets in the statement.

If it does not contain \(2\), it contains \(4\).
It must additionally hit \(\{1,3\}\), \(\{3,5\}\)
and \(\{1,6\}\); the requirement \(\{1,3,5\}\)
is then automatic. If \(1\) is chosen, one of \(3,5\)
is needed. If \(1\) is not chosen, both \(3,6\)
are needed. This gives the last three sets.
Each listed set hits all original members. The
following private witnesses are listed in the
increasing order of the elements of \(H\):
\[
\begin{array}{c|l}
H&(K_c)_{c\in H}\\ \hline
\{1,2\}&\{1,6\},\{2,4\}\\
\{2,3,6\}&\{2,4\},\{1,3,5\},\{1,6\}\\
\{2,5,6\}&\{2,4\},\{1,3,5\},\{1,6\}\\
\{1,3,4\}&\{1,6\},\{2,3,5\},\{2,4\}\\
\{1,4,5\}&\{1,2,3\},\{2,4\},\{2,3,5\}\\
\{3,4,6\}&\{1,2,3\},\{2,4\},\{1,6\}.
\end{array}
\]
Each displayed \(K_c\) is an original conflict
with \(K_c\cap H=\{c\}\). The private witness
criterion proves minimality in every case.
The case split proves exhaustiveness.

For \(\mathcal G\), the singleton conflict forces
\(b\). This hits every member except possibly
\(\{a,c\}\), so precisely one of \(a,c\) is
additionally needed in a minimal answer.
Before pruning, the path sets \(\{a,b\}\) reached
by the edges \(a,b\) and \(b,a\) are the same
shared node. The later label \(\{b\}\) removes
the root's \(a\)-edge. Its \(b\)-edge still reaches
the node \(\{b\}\), whose label \(\{a,c\}\)
has its \(a\)-edge to the shared success
\(\{a,b\}\). It therefore remains reachable.
The \(c\)-edge similarly reaches \(\{b,c\}\).

If instead the second occurrence of \(\{a,b\}\)
were an unconnected closed copy, deleting the
first branch would remove the only stored success
for that answer. This calculation explains the
need for actual shared-node links in the pruning
rule.

In the third family every hitting set contains
both \(a\) and \(b\), and \(\{a,b\}\) is itself
a hitting set. When the root's label is reduced
from \(\{a,b\}\) to \(\{b\}\), the \(b\)-branch
continues to a node labeled \(\{a\}\), then to
the success \(\{a,b\}\). Since
\(\{a\}\not\subseteq\{b\}\), there is no further
root-edge deletion. Failing to update the root's
stored label would apply a different and invalid
comparison with the obsolete set \(\{a,b\}\).
\end{proof}
\end{annexitem}

\begin{annexitem}{annex:p16.diag-circuits}{thm:p16-diag-circuits}{Complete circuit conflict families}
\begin{proof}
Consider an arbitrary set \(N\) of components required
to be normal in the Boolean adder. If
\(\{X_1,X_2\}\subseteq N\), the equations force
\(x=1\) and \(y=0\), contradicting \(y=1\).
If \(\{X_1,A_2,O_1\}\subseteq N\), they force
\(x=v=1\), hence \(z=1\), contradicting \(z=0\).
Both displayed sets are therefore conflicts.

Conversely suppose \(N\) contains neither.
If \(X_1\notin N\), choose
\((x,u,v,y,z)=(0,0,0,1,0)\). This satisfies every
other normal equation. If \(X_1\in N\), then
\(X_2\notin N\), and at least one of \(A_2,O_1\)
is not in \(N\). Set \(x=y=1,u=z=0\).
When \(A_2\in N\), set \(v=1\); then \(O_1\notin N\).
When \(A_2\notin N\), set \(v=0\), which satisfies
the \(O_1\) equation if that component is normal.
All constraints and observed values are satisfied.
Assign abnormality truth values to be false on \(N\)
and true outside it. This gives a full structure.
Thus a normality set is inconsistent exactly when
it contains one of the two displayed conflicts.
Neither displayed set contains the other, so both
are minimal and there are no others.

Their minimal hitting sets are \(\{X_1\}\),
\(\{X_2,A_2\}\) and \(\{X_2,O_1\}\):
either select their common element \(X_1\), or
select \(X_2\) and one of \(A_2,O_1\).
The corresponding wire assignments, in the displayed
order \((x,u,v,y,z)\), are
\[
       (0,0,0,1,0),\qquad(1,0,0,1,0),
                           \qquad(1,0,1,1,0).
\]
There are \(2^3=8\) supersets of the first conflict
and \(2^2=4\) supersets of the second in the
five-component universe. Their union has four
components and hence \(2\) supersets. Inclusion
and exclusion gives \(8+4-2=10\) total conflicts.

For the multiplier network, normal
\(M_1,M_2,A_1\) force \(6+6=10\), a contradiction.
Normal \(M_1,M_3,A_1,A_2\) force
\(x_1=x_3=6\), \(x_2=4\) by \(A_1\), and
\(x_2=6\) by \(A_2\), also a contradiction.
Now let \(N\) contain neither displayed set.
If \(A_1\notin N\), choose \((x_1,x_2,x_3)=(6,6,6)\).
If \(A_1\in N\) and \(M_1\notin N\), choose
\((4,6,6)\). The remaining case has both
\(A_1,M_1\in N\), so \(M_2\notin N\), and
not both \(M_3,A_2\) belong to \(N\).
Choose \((6,4,8)\) when \(A_2\in N\), in which
case \(M_3\notin N\); choose \((6,4,6)\)
when \(A_2\notin N\). These satisfy every required
normal equation. Thus precisely the supersets of
the two displayed sets are conflicts.

Their common elements \(M_1,A_1\) give the two
singleton hitting sets. Omitting both requires
\(M_2\) from the first conflict and one of
\(M_3,A_2\) from the second. This gives exactly
the two double-fault diagnoses.
Wire assignments for the four diagnoses in statement
order are \((4,6,6),(6,6,6),(6,4,8),(6,4,6)\).
With the indicated abnormality assignment, each is
a structure of the appropriate exact theory.
\end{proof}
\end{annexitem}

\begin{annexitem}{annex:p16.diag-measurements}{thm:p16-diag-measurements}{Measurement and minimality}
\begin{proof}
For an old diagnosis \(\Delta\), \(W^-_\Delta\)
is consistent and equivalent to \(W^=_\Delta\).
It has a structure satisfying \(\pi\) exactly when
it does not entail \(\neg\pi\). If it has such
a structure, \(\Delta\) is a withdrawal set for
\(W'\). Every smaller withdrawal for \(W'\)
would also be a smaller withdrawal for \(W\),
which is impossible. Hence \(\Delta\) remains
a diagnosis precisely in this case.

Let \(D'\) be any diagnosis for \(W'\).
A structure of \(W'^-_{D'}\) is also a structure
of \(W^-_{D'}\), so \(D'\) is an old withdrawal
set and contains some old diagnosis \(D\).
If \(D'\) is not itself old, every old diagnosis
\(D\subseteq D'\) is proper. If any such \(D\)
did not predict \(\neg\pi\), the first paragraph
would make it a new diagnosis also, contradicting
minimality of \(D'\). Thus every such \(D\)
predicts \(\neg\pi\). This also proves the
assertion without assuming that all old diagnoses
decide the measurement.

If none of the old diagnoses predicts \(\neg\pi\),
all survive, and no genuinely new one can occur
by the preceding paragraph. When \(W\) is
inconsistent there are no old or new diagnoses.
When only \(W'\) is inconsistent there are no
new diagnoses; the same implications still hold.

If the empty set was not an old diagnosis, any
old diagnosis contained properly in a new one
has at least one component. Its proper superset
has at least two. No new singleton can therefore
appear. If every old diagnosis decides \(\pi\),
the old singletons that do not predict
\(\neg\pi\) are exactly those predicting \(\pi\),
which gives the stated special case.

To see that an unchanged diagnosis family can
still carry new information, take
\(C=\{c\}\), \(W=\{a_c\}\), and a fresh atom
\(q\). The sole diagnosis \(\{c\}\) predicts
neither \(q\) nor \(\neg q\). Measuring \(q\)
leaves that sole diagnosis in place but adds
\(q\) to all of its predictions.
\end{proof}
\end{annexitem}

\begin{annexitem}{annex:p16.diag-measurement-calculations}{thm:p16-diag-measurement-calculations}{Every updated circuit diagnosis}
\begin{proof}
For \(\{M_1\}\), normal \(M_2,M_3\) force
\(x_2=x_3=6\) and normal \(A_1\) forces
\(x_1=4\). For \(\{A_1\}\), all multiplier
outputs equal \(6\). For either diagnosis
containing \(M_2\), normal \(M_1,A_1\) force
\(x_1=6,x_2=4\). These prove the old
predictions.

Add \(x_2=5\). Then \(\{M_2\}\) is a
conflict. The pairs \(\{M_1,A_1\}\) and
\(\{M_3,A_2\}\) are conflicts as well:
the first would force \(6+5=10\), and the
second \(5+6=12\). If a normality set contains
none of these three conflicts, \(M_2\) is
not required normal. Choose \(x_1=6\) if
\(M_1\) is required normal and \(x_1=5\)
otherwise; choose \(x_3=6\) if \(M_3\)
is required normal and \(x_3=7\) otherwise.
Whenever \(A_1\) is normal, \(M_1\) is not,
so its equation holds. The same reasoning
applies to \(A_2\). These assignments prove
that precisely the three listed sets are
minimal conflicts. Every minimal hitting set
contains \(M_2\), one of \(M_1,A_1\), and
one of \(M_3,A_2\). This gives exactly the
four triples in the statement. Their consistent
assignments are among \((5,5,7),(5,5,6),
(6,5,7),(6,5,6)\), in the corresponding order.

Instead add \(x_1=6\). Normal \(A_1\) now
forces \(x_2=4\), so \(\{M_2,A_1\}\) is
a conflict. Combining \(A_1,A_2,M_3\)
also conflicts, since \(x_2=4\) and
\(x_3=6\) cannot sum to \(12\).
To show these are all minimal conflicts,
take a normality set avoiding both.
If \(A_1\) is not required normal, set
\((x_1,x_2,x_3)=(6,6,6)\).
Otherwise \(M_2\) is not normal and at
least one of \(M_3,A_2\) is not normal.
Use \((6,4,8)\) if \(A_2\) is normal
and \((6,4,6)\) otherwise.
These are consistent assignments for every
avoiding normality set. The minimal hitting
sets are therefore \(\{A_1\}\),
\(\{M_2,M_3\}\) and \(\{M_2,A_2\}\),
as claimed.
\end{proof}
\end{annexitem}

\begin{annexitem}{annex:p16.diag-defaults}{thm:p16-diag-defaults}{Maximal consistent defaults and diagnoses}
\begin{proof}
First prove the stated general prerequisite-free
result. Let \(B=\{\alpha_i\}\) be the finite
set of consequent sentences. For a closed candidate
\(E\) the licensed defaults give exactly
\[
 \Gamma(E)=\operatorname{Th}
       (W\cup\{\alpha\in B:\neg\alpha\notin E\}).
\]
This follows also directly from
\cref{def:p18-default-closure}: all prerequisites
are the truth sentence, and no later stage can
license a further default beyond the fixed
candidate test. Write the set inside braces as
\(A_E\).

Suppose \(W\) is consistent and \(E\) is an
extension. It cannot be inconsistent: an
inconsistent closed \(E\) is the whole sentence
language, for which \(A_E=\varnothing\) and
\(\Gamma(E)=\operatorname{Th}(W)\) is consistent.
Thus \(W\cup A_E\) is consistent.
For any \(\alpha\in B\setminus A_E\),
\(\neg\alpha\in E=\operatorname{Th}(W\cup A_E)\),
so adjoining \(\alpha\) would be inconsistent.
This proves maximality of \(A_E\).

Conversely let \(A\subseteq B\) be maximal
consistent with \(W\), and set
\(E=\operatorname{Th}(W\cup A)\).
For \(\alpha\in A\), consistency gives
\(\neg\alpha\notin E\). For
\(\alpha\in B\setminus A\), maximality makes
\(W\cup A\cup\{\alpha\}\) inconsistent,
so classical consequence gives
\(\neg\alpha\in E\). Hence \(A_E=A\),
and \(\Gamma(E)=E\). This proves the general
characterization with all required consistency
conditions.

Apply it to \(B=\{\neg a_c:c\in C\}\).
The maximal jointly consistent normality sets
are exactly \(C\setminus\Delta\) for minimal
withdrawals \(\Delta\), hence for diagnoses.
Even if some designated abnormality sentences
are logically equivalent, maximality includes
every component normality sentence consistent
with the resulting theory, so the index-set
description remains valid.
For a diagnosis, every \(c\in\Delta\) satisfies
\(W^-_\Delta\entails a_c\), while every
\(c\notin\Delta\) has \(\neg a_c\) in that
consistent theory. Thus its extension determines
precisely \(\Delta=\{c:a_c\in E\}\), proving
injectivity and the displayed inverse. Surjectivity
follows from the general maximality description.
The definition of prediction identifies the
extension with all predicted sentences.

If \(W\) is inconsistent, every value of
\(\Gamma\) is the whole sentence language.
That language is consequently its sole fixed
point. Meanwhile no exact or withdrawal theory
has a structure, so there is no diagnosis.
The exceptional case cannot be included in a
bijection between the two families.
\end{proof}
\end{annexitem}

\begin{annexitem}{annex:p16.diag-set-cover}{thm:p16-diag-set-cover}{The complete covering representation}
\begin{proof}
For a normality set \(N\subseteq D\), if
\(R_m\subseteq N\) for an observed \(m\),
the normality implications force every
\(p_d\), \(d\in R_m\), false, while \(o_m\)
and its clause require at least one true.
Thus every such \(R_m\) is a conflict,
including the empty one.

Conversely, suppose no observed \(R_m\) is
contained in \(N\). Assign \(o_m\) true
exactly for \(m\in M^+\), and assign both
\(a_d,p_d\) false for \(d\in N\) and true
for \(d\notin N\). Every normality implication
is satisfied. For each true \(o_m\), the
assumption supplies a \(d\in R_m\setminus N\),
whose \(p_d\) is true. Clauses for unobserved
manifestations have false antecedents.
This is a model of all the required sentences.
Hence conflicts are exactly normality sets
containing some observed \(R_m\), and their
minimal members are exactly the minimal
observed \(R_m\)'s.

With \(N=D\setminus\Delta\), the same criterion
says that \(W^-_\Delta\) is consistent exactly
when each observed \(R_m\) intersects \(\Delta\),
which is the cover condition. Taking minimal
sets proves the diagnosis claim. If
\(M^+=\varnothing\), the construction works
with \(N=D\) and no fault, giving the sole
empty diagnosis. If an observed \(R_m\)
is empty, its forced false consequent makes
the entire base inconsistent and there is
no cover or diagnosis.

The finite propositional representation also
realizes the named first-order presentation:
interpret the disorder and manifestation names
as distinct tagged domain elements, add one
extra element if needed to keep the domain
nonempty, and assign the disorder predicate true exactly
on the named disorders. Set the abnormality
and presence predicates there to the displayed
values, false elsewhere, and set the observation
predicate true exactly on the named observed
manifestations. Conversely every
first-order model supplies those values on
the named constants. The only universal
normality axiom has precisely the required
effect on the named disorders. Thus the
ground construction supplies a full model,
including empty name sets.

A minimum-cardinality cover cannot contain
a smaller cover, so it is inclusion minimal.
For the example, choosing \(a\) covers both
observations. Omitting \(a\) requires both
\(b,c\). These two sets are minimal, and
only the first has cardinality one. This
proves the strict distinction between the
two preference criteria.
\end{proof}
\end{annexitem}

\begin{annexitem}{annex:p16.diag-abduction-boundaries}{thm:p16-diag-abduction-boundaries}{The observation remains a premise of consistency diagnosis}
\begin{proof}
For the first system, \(o\) and normality
\(\neg a_c\) contradict the implication
\(\neg a_c\to\neg o\). Thus the empty set
is not a withdrawal. The valuation
\(a_c=o=\mathbf T\) shows that \(\{c\}\)
is a diagnosis. However, the valuation
\(a_c=\mathbf T,o=\mathbf F\) satisfies
the system and \(a_c\) while falsifying
the observation. This is a countermodel
to the explanatory entailment.

Adding \(a_c\to o\) preserves the first
valuation and the failure of the all-normal
case, so the diagnosis is unchanged.
The added implication and \(a_c\) now
entail \(o\) by the classical truth clause.

In the third system, exactly one of
\(a_c,a_d\) must be true. The two possible
truth assignments therefore realize exactly
the singleton fault sets. Neither truth
assignment is selected by \(W\), since the
other is still a model. Dropping both
normality assumptions leaves \(W\) consistent,
but explicitly setting both abnormalities
true violates its second conjunct.
This proves all three claims and their
distinction from observation-entailing abduction.
\end{proof}
\end{annexitem}

\begin{annexitem}{annex:p16.diag-probability-boundary}{thm:p16-diag-probability-boundary}{Prior fault odds and posterior likelihoods}
\begin{proof}
Independence of the finitely many indicators
gives the product for the exact event that
precisely \(\Delta\) is faulty. Factor out
\(\prod_{c\in C}(1-p_c)>0\) to obtain the
odds product. Adding a component \(c\)
multiplies the probability by
\(p_c/(1-p_c)\), less than one exactly
when \(p_c<1/2\).

For the stated unequal probabilities,
\[
 P(\{a\})=\frac{81}{10^6},
       \qquad P(\{b,c\})=\frac{9999}{10^6}.
\]
Both other components are required normal
in the first exact assignment; \(a\) is
required normal in the second. This verifies
the comparison for complete assignments,
not merely for events saying that some
named faults occur. Independence therefore
does not rank incomparable sets by their
cardinality. If all \(p_c=p<1/2\),
the exact probability is
\(p^{|\Delta|}(1-p)^{|C|-|\Delta|}\),
whose ratio between successive cardinalities
is the common number \(p/(1-p)<1\).

For a positive-probability observed event
\(E\), finite conditional probability gives
\[
 P(\Delta\mid E)=
       \frac{P(E\mid\Delta)P(\Delta)}{P(E)}.
\]
Here \(P(\Delta)>0\) by the assumed
strictly interior fault probabilities.
The denominator is common to all candidate
assignments. A logical model specifies
which assignments and observations can
coexist; it does not assign the numerical
conditional probabilities in this expression.
\end{proof}
\end{annexitem}
